\subsection{The completed characteristic-two cyclic-seven search}

\begin{claim}{B.82}{Proved}{No supersingular factor in the bounded rational-branch family}
Let $H/\overline{\mathbf F}_2$ be a smooth proper connected curve
with a cyclic degree-seven map to the projective line. Suppose its
branch locus is contained in $\mathbf P^1(\mathbf F _8)$ and has at
most seven points. Then $J_H$ has no positive-dimensional
supersingular isogeny factor, even geometrically. Consequently $H$
has no nonconstant map to a supersingular curve of positive genus.
For the normalized branch patterns with four through seven branch
points, the exhaustive counts are respectively
$620,2775,6666,6665$, and every one fails the supersingular-factor
test. This settles precisely the finite search family, not B5.
\end{claim}

\begin{proof}
Put $\mathbf F _8=\mathbf F _2[\alpha]/(\alpha^3+\alpha+1)$.
There are at least two unbranched rational points on the projective
line. A projective linear transformation over $\mathbf F _8$ can
send two of them to $0,\infty$ and one branch point to $1$.
Kummer theory, since seven is prime to two, and a change of cyclic
generator give the geometric model
\[
 y^7=\prod_{i=0}^6(x-\alpha^i)^{e_i},\qquad
 e_0=1,\quad e_i\in\{0,1,\ldots,6\},\quad
                       \sum_i e_i=0\pmod7.                 \tag{FC.7}
\]
Indeed valuations give the exponents modulo seven, and a rational
function with all valuations divisible by seven is a constant times
a seventh power on the projective line. A constant becomes a seventh
power over the algebraic closure. Thus this normalization loses no
geometric isomorphism class in the stated family. Let $r$ be the
number of nonzero exponents. The cases $r\le2$ have genus zero or
are impossible; the case $r=3$ is treated below without a search.

Let $\zeta$ be a primitive seventh root of unity and let
$\chi_m:\mathbf F _{8^m}^{\times}\to\langle\zeta\rangle$
be the norm-compatible character with $\chi_1(\alpha)=\zeta$.
At an unramified finite point use $\chi_m(f(x))$, at a ramified
point use zero, and at infinity use one, because the degree in
(FC.7) is divisible by seven and its leading coefficient is one.
Write $S_m$ for the sum of these values. The nontrivial character
factor is
\[
 L(T)=\exp\left(\sum_{m\ge1}S_mT^m/m\right)
     =\sum_{j=0}^{r-2}a_jT^j,
 \qquad a_0=1,\qquad
             j a_j=\sum_{m=1}^jS_ma_{j-m}.                 \tag{FC.8}
\]
Here is a direct degree bound. The finite Euler product is the
Dirichlet $L$-series of the character obtained by evaluating monic
polynomials at the $r$ branch points and taking the prescribed
powers of the residue characters. For degree at least $r$, the
values at these points range uniformly and independently, so the
coefficient sum vanishes by the nontrivial character sum in any one
coordinate. The character is trivial on constants because
$\sum e_i=0$ modulo seven. Its finite $L$-polynomial therefore has
the factor $1-T$: sum over all polynomials of degree below $r$,
first grouping by their nonzero scalar multiple, to obtain zero
again by uniform evaluation. Restoring the unramified Euler factor
at infinity divides by $1-T$. This proves $\deg L\le r-2$.
The six nontrivial factors are its conjugates under
$\zeta\mapsto\zeta^j$, $1\le j\le6$; their product is the zeta
numerator of $H$, since the invariant factor is the zeta function
of the projective line. These statements also follow directly by
multiplying the seven local Kummer Euler factors. Thus the finite
sums through $m=r-2$ determine every root that needs testing.

We next prove that the test list is complete. A supersingular
Frobenius root $\lambda$ over $\mathbf F _8$ has
$\lambda/\sqrt8$ a root of unity. To justify this implication,
all its finite valuations, after this division, are zero: the
Newton slopes at two are $1/2$, and the other valuations vanish
because a Weil number has norm a power of two. All complex
absolute values are one. An algebraic unit with this property is
a root of unity: the monic polynomials of its powers have bounded
integer coefficients and bounded degree, so only finitely many
powers can occur among their roots. This is the needed elementary
form of Kronecker's argument.

Put $K=\mathbf Q(\zeta)$. If $\lambda$ is a root of one of our
character factors, then $[K(\lambda):K]\le5$. The root of unity
$\beta=\lambda^2/8$ has order $N$, whence
\[
 [K(\beta):K]=\frac{\varphi(\operatorname{lcm}(7,N))}{6}\le5.
\]
Write $N=7^a m$ with $7\nmid m$. Necessarily $a\le1$ and
$\varphi(m)\le5$, so
\[
                 m\in\{1,2,3,4,5,6,8,10,12\}.
\]
After multiplying $\lambda$ by a power of $\zeta$, its square
divided by eight has order $m$. The resulting number belongs to
a cyclotomic field ramified only at primes dividing $2m$. That
field is linearly disjoint from $K$: every nontrivial subfield of
$K$ is ramified at seven. Its degree over $K$ is therefore its
degree over $\mathbf Q$.

Factoring
$\operatorname{Res}_Z(\Phi_m(Z),U^2-8Z)$ for the nine displayed
values of $m$ now gives exactly the following possible minimal
polynomials of degree at most five:
\[
\begin{gathered}
 U^2-8,\quad U^2+8,\quad U^2\pm4U+8,\\
 U^4\pm8U^2+64,\qquad
 U^4\pm4U^3+8U^2\pm32U+64,
\end{gathered}                                                \tag{FC.9}
\]
where the two signs in the last expression agree. The remaining
factors, for $m=5,8,10$, have degree eight and are irreducible.
These small rational factorizations, with every coefficient, are
also checked symbolically in the certificate. The listed factors
are irreducible over $K$ by the same disjointness argument.
They all have degree at least two, which proves the assertion for
$r=3$ immediately.

For each of the $16726$ normalized exponent vectors with $r\ge4$,
the program computes (FC.8) in the integral basis
$1,\zeta,\ldots,\zeta^5$ and tests divisibility by the reciprocal
polynomial of each of (FC.9), after all seven changes
$T\mapsto\zeta^jT$. No division succeeds. The complete character
polynomials and every power-sum residue count are retained, not
merely the absence of a search hit. The exact arithmetic uses the
fields of degrees $3,6,9,12,15$ over $\mathbf F _2$ with respective
binary moduli $11,67,529,4179,32771$; the full nonzero power table
checks that each proposed modulus gives a field and each chosen
generator has full order. The embeddings of $\alpha$ and norm
characters are recorded as well. This exhausts (FC.7) and (FC.9),
so no Frobenius root is supersingular.

The conclusion persists under every finite constant extension:
if $\lambda^n/8^{n/2}$ were a root of unity, then
$\lambda/\sqrt8$ would be one. Finally a nonconstant map to a
positive-genus supersingular curve induces pullback and norm on
Jacobians whose composition is multiplication by the degree.
Its Jacobian would therefore be a positive-dimensional
supersingular isogeny factor, a contradiction. This argument also
covers inseparable morphisms.
\end{proof}

\paragraph{Check record and exact boundary.}
Checked only in this writing: GPT-6 Astra Pro, finite-checks reader,
18 November 2026. We completed the supplied search, proved completeness
of the previously heuristic eight-factor list, and retained all
$16726$ character polynomials. The program
\path{code/b5_cyclic7_complete.py} is an inspected adaptation of the
supplied search, not a blind independent implementation. Its full
output and compressed polynomial certificate are retained. A separate verifier, \path{code/b5_cyclic7_verify.py}, imports no
producer code. It checks all Newton identities by cyclic convolution,
then proves nondivisibility by descending polynomial division after
specializing $\zeta$ at split primes. Of its $841596$ nondivisibility
checks, $840805$ have a nonzero remainder modulo $29$ and the remaining
$791$ have one modulo $43$. The leading coefficients are units in
these fields, so each nonzero remainder is a valid characteristic-zero
nondivisibility certificate. The effort
setting is not separately exposed. Branch sets with eight or nine
$\mathbf F _8$-points, branch points over larger fields, and arbitrary
prime-to-two correspondences are outside this statement. In particular
this negative finite family neither proves nor refutes B5.
